% =============================================================
% 第4章 技術成熟度レベルTRLの推移 (課題7) 章本体
% reports.tex から \input される．
% =============================================================

\section{はじめに}
本章には，Arduinoオーケストラ開発における技術成熟度レベル（TRL: Technology Readiness Level）の評価を，
FBS（機能分解構造）の各機能ごとに記録する．
TRLは NASA 由来の評価指標で，1（基礎原理確認）から9（実運用実証済み）までの9段階で技術成熟度を示す．
本プロジェクトでは，FBSで定義された各機能要素に対して個別にTRLを評価することで，
機能ごとの開発進捗・成熟度を可視化する．

\section{TRLレベルの定義（参考）}

\begin{table}[H]
  \centering
  \begin{tabularx}{\linewidth}{cX}
    \toprule
    レベル & 概要 \\
    \midrule
    TRL 1 & 基礎原理が観察・報告された \\
    TRL 2 & 技術コンセプトおよび応用が定式化された \\
    TRL 3 & 解析・実験により概念が実証された \\
    TRL 4 & 構成要素が試験環境で検証された \\
    TRL 5 & 構成要素が運用に近い環境で検証された \\
    TRL 6 & システム試作モデルが運用に近い環境で実証された \\
    TRL 7 & システム試作モデルが運用環境で実証された \\
    TRL 8 & 実システムが完成し，試験で適格性を確認した \\
    TRL 9 & 実システムが実運用で実証された \\
    \bottomrule
  \end{tabularx}
\end{table}

\section{FBS要素ごとのTRL評価}

FBSで定義された各機能要素の現時点におけるTRLレベルと，その評価根拠を表に示す．
FBSの構造はマネジメント計画書（\texttt{management\_plan.tex}）の機能分解構造節を参照．

\begingroup
\small
\begin{longtable}{p{3cm}p{4.5cm}p{1.8cm}p{6cm}}
\toprule
FBSフェーズ & 機能 & TRLレベル & 根拠 \\
\midrule
\endfirsthead

\toprule
FBSフェーズ & 機能 & TRLレベル & 根拠 \\
\midrule
\endhead

\bottomrule
\endfoot

% --- 指揮動作検出 ---
指揮動作検出 & 拍タイミング検出 &  &  \\
 & テンポ推定 &  &  \\
 & フェルマータ検出 &  &  \\
 & 小節同期見逃し補完 &  &  \\
\midrule

% --- 指揮者人形動作 ---
指揮者人形動作 & 腕往復動作 &  &  \\
 & 輝度エンコードLED制御 &  &  \\
\midrule

% --- 演奏制御 ---
演奏制御 & 輪唱オフセット制御 &  &  \\
 & 楽譜進行制御 &  &  \\
 & テンポ追従演奏 &  &  \\
\midrule

% --- 音響合成 ---
音響合成 & 音色合成（倍音・ADSR） &  &  \\
 & エフェクト付与 &  &  \\
 & フェルマータ延奏 &  &  \\
\midrule

% --- 状態管理 ---
状態管理 & 演奏状態遷移管理 &  &  \\
 & 異常時フォールバック &  &  \\

\end{longtable}
\endgroup
